% -----------------------------------------------------------------------
% Packages
% -----------------------------------------------------------------------
\usepackage[utf8]{inputenc}
\usepackage[T1]{fontenc}
\usepackage{amsmath, amssymb}
\usepackage{graphicx}
\usepackage{booktabs}
\usepackage{threeparttable}
\usepackage{array}
\usepackage{hyperref}
\usepackage[english]{babel}
\usepackage[useregional]{datetime2}
\usepackage{cleveref}
\usepackage{xstring}
\usepackage{geometry}
\usepackage{mdframed}
\usepackage{tcolorbox}       % styled coloured boxes
  \tcbuselibrary{skins, breakable}
\usepackage{xcolor}
\usepackage{soul}
\usepackage{setspace}
\usepackage{needspace}
\usepackage[nopatch=footnote]{microtype} % CHANGED vs src (was: \usepackage{microtype})
\usepackage[shortlabels]{enumitem}
% Fonts
\usepackage{helvet}
\usepackage{times}
\usepackage{titlesec}   % sections titles
\usepackage{eurosym}
\usepackage{changepage}
\usepackage{wasysym}
\usepackage{fancyhdr}
\usepackage{lastpage}

% Line numbering (for editorial review)
% -------------------------------------------------------------
\usepackage{lineno}
\renewcommand{\linenumberfont}{\normalfont\tiny\sffamily\color{gray}}
\AtBeginDocument{\linenumbers}


% Colours
\definecolor{headingcolor}{RGB}{42, 110, 184}
\definecolor{boxcolor}{RGB}{244, 247, 252}

% -----------------------------------------------------------------------
% Part-dependent box colour
% -----------------------------------------------------------------------
% Part 1: Chapters 0-5   -> light blue
% Part 2: Chapters 6-9   -> light green
% Part 3: Chapters 10-15 -> light amber
\definecolor{boxcolorpart1}{RGB}{244, 247, 252}
\definecolor{boxcolorpart2}{RGB}{238, 252, 250}
\definecolor{boxcolorpart3}{RGB}{252, 246, 238}

\newcommand{\boxcolor}{%
  \ifnum\value{chapter}<6
    \colorlet{boxcolorcurrent}{boxcolorpart1}%
  \else\ifnum\value{chapter}<10
    \colorlet{boxcolorcurrent}{boxcolorpart2}%
  \else
    \colorlet{boxcolorcurrent}{boxcolorpart3}%
  \fi\fi
}


\usepackage[font={footnotesize,sf}, labelfont={bf,color=headingcolor}]{caption}

\newcommand{\degree}{\ensuremath{{}^{\circ}}}

\hypersetup{
  colorlinks  = true,        % coloured text instead of rectangles
  linkcolor   = blue!60!black,  % internal links (ToC, \ref, \pageref)
  citecolor   = blue!60!black,  % bibliography citations
  urlcolor    = blue!60!black,  % external URLs
}

% Page layout
\geometry{
 a4paper,
 left=40mm,
 right=20mm,
 top=20mm,
 bottom=25mm
}

% Text body tab
\newcommand{\bodytab}{10mm}

% Paragraph indentation
\setlength{\parindent}{0mm}

% Paragraph spacing
\setlength{\parskip}{2.5mm}

% --------------------------------------------------------------
% Line spacing
\linespread{1.1}

% --------------------------------------------------------------
% Page style
\pagestyle{fancy}

% ----------- header -----------
\fancyhead[L]{\sffamily\scriptsize{Sivapalan \textit{et al.} (in press)}}
\fancyhead[R]{\sffamily\scriptsize{\chaptername~\thechapter\;--\;Manuscript}}
\renewcommand{\headrulewidth}{0.2pt}

% ----------- footer -----------
\fancyfoot[L]{\sffamily\scriptsize{Compilation:~\DTMNow}}
\fancyfoot[C]{\sffamily\scriptsize{}}
\fancyfoot[R]{\sffamily\scriptsize{\thepage\;/\;\pageref{LastPage}}}
\renewcommand{\footrulewidth}{0.2pt}

% Set sans-serif font
\renewcommand{\sfdefault}{phv} % phv=Helvetica

% Chapter and section title formatting
\titleformat{\chapter}[block]
  {\bfseries\fontfamily{phv}\selectfont\Large\color{headingcolor}}
  {\chaptername~\thechapter}{1em}{}
\titleformat*{\section}{\bfseries\fontfamily{phv}\selectfont\large\color{headingcolor}}
\titleformat*{\subsection}{\bfseries\fontfamily{phv}\selectfont\normalsize\color{headingcolor}}
\titleformat*{\subsubsection}{\bfseries\fontfamily{phv}\selectfont\small\color{headingcolor}}
\setcounter{secnumdepth}{3}
\renewcommand{\thesubsubsection}{(\alph{subsubsection})}
% -----------------------------------------------------------------------
% Standard commands
% -----------------------------------------------------------------------

% Author/editorial note
\newcommand{\authornote}[1]{\hl{[Author note: #1]}}

% -----------------------------------------------------------------------
% Figure tier resolution (T1 first, then T0)
% -----------------------------------------------------------------------
% \figlocate{<figure key>}  e.g. \figlocate{C07-mtx-compartments}
% Sets:
%   \figchapnum  chapter number taken from characters 2-3 of the key (e.g. 07)
%   \figtier     T1 or T0, whichever was selected
%   \figbase     path WITHOUT extension, e.g. ./figures/chapter07/T1/<key>_T1
% Search order (first hit wins):
%   1. figures/chapterNN/T1/<key>_T1.jpeg
%   2. figures/chapterNN/T1/<key>_T1.png
%   3. figures/chapterNN/T0/<key>_T0.jpeg
%   4. figures/chapterNN/T0/<key>_T0.png
% If nothing is found, \figbase and \figtier fall on T0 and the macros draw
% their empty placeholder frame, exactly as before.
% The macros that call \figlocate keep their own \IfFileExists{\figbase.jpeg}
% / \IfFileExists{\figbase.png} tests, so jpeg-before-png and the per-format
% layouts are unchanged.
\newcommand{\figlocate}[1]{%
  \StrLeft{#1}{3}[\figtmpA]%
  \StrRight{\figtmpA}{2}[\figchapnum]%
  \IfFileExists{./figures/chapter\figchapnum/T1/#1_T1.jpeg}{%
    \def\figtier{T1}%
  }{%
    \IfFileExists{./figures/chapter\figchapnum/T1/#1_T1.png}{%
      \def\figtier{T1}%
    }{%
      \def\figtier{T0}%
    }%
  }%
  \edef\figbase{./figures/chapter\figchapnum/\figtier/#1_\figtier}%
}

% -----------------------------------------------------------------------
% Opening image placeholder (Tier 1 — no float, no numbering)
% #1 = caption text
\newcommand{\openingimage}[2]{%
\figlocate{#1}%
\begin{adjustwidth}{0cm}{0cm}
    \centering
    \IfFileExists{\figbase.jpeg}{%
      \parbox[c][4cm][c]{4cm}{\centering\sffamily\footnotesize Figure \\ #1}%
      \hspace{0.5cm}%
      \parbox[c][6cm][c]{9cm}{\centering\includegraphics[height=6cm, width=9cm,keepaspectratio]{\figbase.jpeg}}%
    }{%
      \IfFileExists{\figbase.png}{%
        \parbox[c][4cm][c]{4cm}{\centering\sffamily\footnotesize Figure \\ #1}%
        \hspace{0.5cm}%
        \parbox[c][6cm][c]{12cm}{\centering\includegraphics[height=6cm, width=12cm,keepaspectratio]{\figbase.png}}%
      }{%
        \fbox{%
          \parbox{0.7\linewidth}{%
            \centering
            \vspace{2cm}
            \sffamily \footnotesize Figure #1
            \vspace{2cm}
          }%
        }%
      }%
    }
  \end{adjustwidth}
  {\sffamily \footnotesize #2}
\vspace{5mm}
}

\newcommand{\inthischapter}[1]{
\noindent \textit{In this chapter} \ldots

\noindent #1

}

% -----------------------------------------------------------------------
% Didactic learning checks
% Uses the same blue as chapter/section titles and the normal body (Times) font.
% Label and prompt remain in one paragraph; vertical spacing is added before/after.
\newcommand{\didacticitem}[2]{%
  \par\Needspace{5\baselineskip}\addvspace{3mm}%
  \noindent{\normalfont\color{headingcolor}\textbf{#1:}\ #2}%
  \par\addvspace{3mm}%
}
\newcommand{\howtousechapter}[1]{\didacticitem{How to use this chapter}{#1}}
\newcommand{\learningcheck}[2]{\didacticitem{Learning check -- #1}{#2}}
\newcommand{\finalselfcheck}[1]{\didacticitem{Final self-check}{#1}}

% -----------------------------------------------------------------------
% BOXES
% Shared style preset — edit once, applies to both boxes below
\tcbset{
  myboxstyle/.style={
    enhanced,              % full TikZ engine; needed for most options
    breakable,             % allow splitting across pages
    colback      = boxcolorcurrent,
    colbacktitle = boxcolorcurrent,
    coltitle     = headingcolor,           % white title text
    colframe     = headingcolor,   % frame line same as title bar
    boxrule      = 1pt,           % thin frame line; use 0pt + "frame hidden" for none
    arc          = 0pt,             % slightly rounded corners; arc=0pt for sharp
    left         = 1mm,  % left inner padding
    right        = 1mm,  % right inner padding
    fonttitle    = \bfseries\sffamily\small,
    fontupper    = \sffamily\footnotesize,
    fontlower    = \sffamily\footnotesize,
  }
}

% Numbered variant — counter increments automatically.
% Takes an optional first argument for per-instance tcolorbox keys,
% most usefully: \begin{keybox}[list text=Short entry]{Full title}
% The optional arg is passed directly to tcolorbox so any valid key works.
\newtcolorbox[
  auto counter,          % create and manage a counter for this box type
  number within = chapter,  % reset counter at each new chapter
  list inside   = lob,   % register entries for a \listof if desired
]{boxmain}[3][]{          % #1 = optional keys, #2 = theme #3 = title
  myboxstyle,
  title = {Box~\thetcbcounter\;#2: #3},
  #1,                    % per-instance keys applied last (override style)
}


% Professional Skills
\newtcolorbox{boxskills}{
  myboxstyle,
  title = {Professional practice relevance and skills developed},
}

% Biographies (portrait on the left, text on the right)
\newtcolorbox{boxbio}[3]{
  myboxstyle,
  breakable=false,
  sidebyside,
  sidebyside align=top,
  lefthand width=0.2\linewidth,
  title = {Biography: #1 (#2 -- #3)},
}

% Portrait placeholder for boxbio (Tier 1 -- resolved at Tier 2)
% #1 = caption text shown in the placeholder frame (typically the person's name)
\newcommand{\boxbiofigplaceholder}[1]{%
  \figlocate{#1}%
  \centering
  \IfFileExists{\figbase.jpeg}{%
    \parbox[c][4cm][c]{3cm}{\centering\includegraphics[height=4cm, width=3cm,keepaspectratio]{\figbase.jpeg}}%
    \newline
    \parbox[c][1cm][c]{3cm}{\centering\sffamily\footnotesize Figure \\ #1}%
    }{%
      \IfFileExists{\figbase.png}{%
        \parbox[c][4cm][c]{3cm}{\centering\includegraphics[height=4cm, width=3cm,keepaspectratio]{\figbase.png}}%
        \newline
        \parbox[c][1cm][c]{3cm}{\centering\sffamily\footnotesize Figure \\ #1}%
      }{%
        \fbox{\parbox[c][4cm][c]{3cm}{\centering \sffamily \footnotesize Figure \\ #1}}%
      }%
    }%
}%

% Take-home messages
\newtcolorbox{boxmessages}{
  myboxstyle,
  title = {Take home messages of Chapter \thechapter},
}

% Make \cref / \Cref recognise the boxmain counter
\crefname{tcb@cnt@boxmain}{box}{boxes}
\Crefname{tcb@cnt@boxmain}{Box}{Boxes}


% --------------------------------------------------------------------
% Figure placeholder
% #1 = positional key  e.g. C07-mtx-compartments
% #2 = display label   e.g. C07-mtx-compartments
% #3 = caption text
\newcommand{\figplaceholder}[3]{%
  \figlocate{#1}%
  \begin{figure}[htbp]
  \begin{adjustwidth}{0cm}{0cm}
    \centering
    \IfFileExists{\figbase.jpeg}{%
      \parbox[c][4cm][c]{4cm}{\centering\sffamily\footnotesize Figure \\ #2}%
      \hspace{0.5cm}%
      \fbox{\parbox[c][4cm][c]{10cm}{\centering\includegraphics[height=4cm, width=10cm,keepaspectratio]{\figbase.jpeg}}}%
    }{%
      \IfFileExists{\figbase.png}{%
        \parbox[c][4cm][c]{4cm}{\centering\sffamily\footnotesize Figure \\ #2}%
        \hspace{0.5cm}%
        \fbox{\parbox[c][4cm][c]{10cm}{\centering\includegraphics[height=4cm, width=10cm,keepaspectratio]{\figbase.png}}}%
      }{%
        \fbox{%
          \parbox{0.7\linewidth}{%
            \centering
            \vspace{2cm}
            \sffamily \footnotesize Figure #2
            \vspace{2cm}
          }%
        }%
      }%
    }
    \caption{#3}
    \label{#1}
  \end{adjustwidth}
  \end{figure}%
}

\newcommand{\boxfigplaceholder}[1]{%
\figlocate{#1}%
\begin{adjustwidth}{0cm}{0cm}
    \centering
    \IfFileExists{\figbase.jpeg}{%
      \parbox[c][4cm][c]{4cm}{\centering\sffamily\footnotesize Figure \\ #1}%
      \hspace{0.5cm}%
      \fbox{\parbox[c][4cm][c]{8cm}{\centering\includegraphics[height=4cm, width=8cm,keepaspectratio]{\figbase.jpeg}}}%
    }{%
      \IfFileExists{\figbase.png}{%
        \parbox[c][4cm][c]{4cm}{\centering\sffamily\footnotesize Figure \\ #1}%
        \hspace{0.5cm}%
        \fbox{\parbox[c][4cm][c]{8cm}{\centering\includegraphics[height=4cm, width=8cm,keepaspectratio]{\figbase.png}}}%
      }{%
        \fbox{%
          \parbox{0.7\linewidth}{%
            \centering
            \vspace{2cm}
            \sffamily \footnotesize Figure #1
            \vspace{2cm}
          }%
        }%
      }%
    }
  \end{adjustwidth}
}

% -----------------------------------------------------------------------
% Environments
% -----------------------------------------------------------------------


% Worked examples (theorem-like, numbered automatically per chapter)
% -----------------------------------------------------------------------
\newcounter{workedexample}
\numberwithin{workedexample}{chapter}
\renewcommand{\theworkedexample}{\thechapter.\arabic{workedexample}}

\newenvironment{workedexample}[1]
{%
  \refstepcounter{workedexample}%
  \setlength{\parindent}{0pt}%
  \par\medskip\noindent
  {\bfseries\small\fontfamily{phv}\selectfont\color{headingcolor}%
   Worked Example~\theworkedexample: #1}\par
}
{%
  \par\hfill End of Worked Example $\blacksquare$\par
}

% Problem / Solution substructure within a worked example
\newcommand{\subcase}[1]{\par\noindent\textbf{#1.}\enspace\ignorespaces}
% Make \cref / \Cref recognise the workedexample counter
\crefname{workedexample}{Worked Example}{Worked Examples}
\Crefname{workedexample}{Worked Example}{Worked Examples}


% -------------------------------------------------------------
% prefix numbering
\numberwithin{equation}{chapter}
\numberwithin{figure}{chapter}
\numberwithin{table}{chapter}


% ======================================================================
% Homework assimilation support (auto-generated by assimilate_homework.py)
% Source preamble: Hydrology_Homework_Material_Chapters_1-15_TeXworks.tex
% ======================================================================
\usepackage{longtable}
\colorlet{titleblue}{headingcolor}
\newenvironment{sourceanchor}{%
  \par\smallskip\begingroup\small\itshape\color{headingcolor}%
}{\par\endgroup\smallskip}
\newenvironment{hangpart}[1]{%
  \par\noindent\hangindent=1.7em\hangafter=1\textbf{(#1)}\ %
}{\par}

% --- addons from preamble_addons.tex ---
% Hand-maintained additions appended to out/config.tex by assimilate_homework.py.
% Keep every entry safe for all chapters (it is shared by the whole book).

% Unicode degree sign: works in text and math mode (the chapters use a literal
% degree sign inside $...$, which inputenc maps to the text-only \textdegree).
\DeclareUnicodeCharacter{00B0}{\ifmmode{}^{\circ}\else\textdegree\fi}